% Generated by scripts/benchmarks_to_latex.py from results/benchmarks.json.
% Do not edit by hand: re-run the script after scripts/run_benchmarks.py.

Ho misurato le prestazioni con \texttt{scripts/run\_benchmarks.py} e ho fissato delle soglie
di regressione in \texttt{pytest -m bench}, così se una modifica peggiora i tempi me ne accorgo.
Salvo dove indicato, ho usato l'intero dataset (3\,730\,870 snapshot,
1{,}19\,GB in \texttt{float64}) già caricato in RAM.

\paragraph{Ambiente di misura.} Il mio portatile: CPU Intel Core Ultra 7 155H
(22 thread OpenMP), Windows 11,
compilatore gcc (MinGW-w64) con \texttt{-O3 -march=native -fopenmp},
Python 3.13.9, NumPy 2.5.3, motore v1.0.0.
Per i tempi riporto il migliore su più ripetizioni; per le latenze delle singole chiamate
riporto i percentili su centinaia di migliaia di chiamate.

\subsection{Calcolo dei segnali sull'intero dataset}

Nella Tabella~\ref{tab:bench-headline} confronto il mio codice originale, che attraversava il
confine Python/C una o due volte per ogni snapshot, con l'API batch, che passa l'intero
array al motore C in una sola chiamata.

\begin{table}[H]
\centering
\caption{Tempo per calcolare i segnali su tutti gli snapshot.}
\label{tab:bench-headline}
\begin{tabular}{p{7.2cm}rrr}
\toprule
Percorso & Tempo totale (s) & ns/snapshot & vs originale \\
\midrule
Script OBI originale (\texttt{csv.reader} + 2 chiamate ctypes per tick) & 24{,}6 & 6\,582 & 1{,}0$\times$ \\
Loop WOBI originale (40 scritture ctypes + 1 chiamata per tick) & 34{,}1 & 9\,132 & 0{,}7$\times$ \\
WOBI vettoriale NumPy (1 thread) & 0{,}972 & 261 & 25$\times$ \\
WOBI batch in C, 1 thread & 0{,}087 & 23{,}4 & 282$\times$ \\
\textbf{WOBI batch in C, 22 thread} & 0{,}028 & 7{,}6 & 872$\times$ \\
\bottomrule
\end{tabular}
\end{table}

Il mio script originale impiega 24{,}6\,s sul CSV completo (nel primo
README avevo scritto circa 22\,s). Profilandolo ho capito che quasi tutto quel tempo era
interprete Python e marshalling ctypes, non calcolo: il C che avevo scritto pesava pochissimo.
Con il batch lo stesso lavoro scende a
28\,ms, un fattore
872.

\subsection{Scalabilità OpenMP}

\begin{table}[H]
\centering
\caption{Kernel batch in C al variare del numero di thread (dati già in RAM).}
\label{tab:bench-scaling}
\small\setlength{\tabcolsep}{4pt}
\begin{tabular}{rrrrrrr}
\toprule
Thread & OBI ns/riga & Speedup & WOBI ns/riga & Mrighe/s & Speedup & Efficienza \\
 & & OBI & & WOBI & WOBI & WOBI \\
\midrule
1 & 23{,}14 & 1{,}0$\times$ & 23{,}35 & 43 & 1{,}0$\times$ & 100\,\% \\
2 & 11{,}13 & 2{,}1$\times$ & 14{,}24 & 70 & 1{,}6$\times$ & 82\,\% \\
4 & 6{,}01 & 3{,}8$\times$ & 9{,}96 & 100 & 2{,}3$\times$ & 59\,\% \\
8 & 4{,}00 & 5{,}8$\times$ & 8{,}89 & 112 & 2{,}6$\times$ & 33\,\% \\
12 & 3{,}73 & 6{,}2$\times$ & 7{,}79 & 128 & 3{,}0$\times$ & 25\,\% \\
16 & 3{,}84 & 6{,}0$\times$ & 7{,}69 & 130 & 3{,}0$\times$ & 19\,\% \\
22 & 3{,}85 & 6{,}0$\times$ & 7{,}55 & 132 & 3{,}1$\times$ & 14\,\% \\
\bottomrule
\end{tabular}
\end{table}

\begin{figure}[H]
\centering
\includegraphics[width=0.8\textwidth]{openmp_scaling.png}
\caption{Speedup dei kernel batch rispetto a un thread. La linea tratteggiata è la
scalabilità lineare ideale.}
\label{fig:bench-scaling}
\end{figure}

La scalabilità si appiattisce presto (Figura~\ref{fig:bench-scaling}): con
22 thread lo speedup è 6{,}0$\times$ per OBI e
3{,}1$\times$ per WOBI. Ogni snapshot occupa 320 byte ma i kernel
eseguono poche operazioni per byte letto (OBI usa 16 byte su 320), quindi il limite è la
banda di memoria e non il numero di core. Il margine residuo sta nel layout dei dati
(colonne separate, precisione ridotta), non in più thread.

\subsection{Latenza per snapshot (tick-to-signal)}

Un sistema live riceve uno snapshot alla volta e deve produrre il segnale subito, quindi ho
misurato anche il costo di una singola chiamata (Tabella~\ref{tab:bench-latency}).

\begin{table}[H]
\centering
\caption{Latenza per singolo snapshot, in nanosecondi (massimo in microsecondi).}
\label{tab:bench-latency}
\begin{tabular}{p{6.2cm}rrrr}
\toprule
Percorso & p50 & p99 & p99{,}9 & max ($\mu$s) \\
\midrule
Solo C, \texttt{lob\_obi\_one} & 16 & 317 & 547 & 1\,699 \\
Solo C, \texttt{lob\_wobi\_one} & 21 & 62 & 81 & 1\,854 \\
Python $\to$ ctypes $\to$ \texttt{lob\_obi\_one} & 300 & 800 & 1\,200 & 1\,853 \\
Python $\to$ ctypes $\to$ \texttt{lob\_wobi\_one} & 1\,000 & 2\,700 & 11\,000 & 2\,931 \\
WOBI originale per tick & 7\,200 & 19\,100 & 94\,802 & 7\,287 \\
\bottomrule
\end{tabular}
\end{table}

Il kernel WOBI in C costa circa 21\,ns (misurati con il
contatore TSC del processore, overhead del timer 17\,ns
incluso). Chiamato da Python la mediana sale a 1\,000\,ns, circa
48 volte di più: il costo è il confine Python/ctypes, non il calcolo.
Per le misure lato Python ho dovuto usare \texttt{perf\_counter\_ns}, che su Windows ha
risoluzione di 100\,ns, quindi quei valori sono quantizzati a quel passo. I massimi, nell'ordine dei millisecondi, sono dovuti
a preemption del sistema operativo e interrupt, non al codice.

\subsection{Caricamento dei dati e pipeline completa}

\begin{table}[H]
\centering
\caption{Caricamento dei dati e pipeline di ricerca end-to-end sull'intero dataset.}
\label{tab:bench-e2e}
\begin{tabular}{lr}
\toprule
Passo & Tempo \\
\midrule
Parsing del CSV grezzo (una tantum, \texttt{prepare\_data.py}) & circa 13\,s \\
Apertura della cache \texttt{.npy} memory-mapped & 1{,}2\,ms \\
Caricamento completo della cache in RAM & 0{,}56\,s \\
\midrule
Segnale OBI (batch C, da memory-map) & 164\,ms \\
WOBI e micro-price (batch C) & 133\,ms \\
Backtest OBI a orizzonte fisso (541\,959 trade) & 0{,}84\,s \\
Backtest WOBI con bracket (62\,921 trade) & 0{,}63\,s \\
\midrule
\textbf{Totale pipeline} & \textbf{1{,}77\,s} \\
\bottomrule
\end{tabular}
\end{table}

Una volta ridotti i segnali a decine di millisecondi, il collo di bottiglia è diventato il
backtest, che itera in Python una volta per trade. Rendendo vettoriale il calcolo di ingressi,
uscite e PnL ho portato il backtest OBI da 5{,}2\,s a meno di 1\,s, con gli stessi trade
selezionati.

\subsection{Sintesi}

\begin{itemize}[leftmargin=*]
\item Il calcolo dei segnali non è più un collo di bottiglia: 872 volte
più veloce del codice originale sull'intero dataset.
\item I kernel batch sono limitati dalla banda di memoria; il prossimo guadagno viene dal
layout dei dati, non da più thread.
\item Nel percorso per singolo tick oltre il 95\,\% del tempo è overhead Python/ctypes;
un sistema live dovrebbe tenere l'intero percorso critico in C/C++.
\end{itemize}
